Jesus sits on the mountain, and his disciples come near to hear him.
He promises them the kingdom of heaven. As he describes their life,
he speaks of mourning, hunger for justice, and persecution. These
people are blessed while they still have something to endure. On the
feast of All Saints, we hear Christ describe the holiness of those we
honor, and the life he is teaching us to share.

Before this Gospel, the Alleluia gives us his invitation:
``Come to me all you that labor and are burdened, and I will refresh
you.'' We need to understand that promise beside the words he speaks
on the mountain. He offers rest to people who will still need courage.
If we expect coming to Christ to remove every difficulty, his blessing
of the persecuted will be hard to understand. His own explanation tells
us what kind of rest he gives.

In the verses following that invitation, Jesus speaks of taking his
yoke and learning from him. A yoke is something worn for work. He asks
us to learn his meekness and humility of heart. Saint John Chrysostom
connects those words with the beginning of the Beatitudes. The humble
person can begin to receive Christ's teaching, and in following it
begin to be freed from burdens that sin has imposed.

Chrysostom gives a plain example. A person who strikes back continues
a quarrel. Refusing to retaliate can bring that exchange to an end.
The restraint costs something: anger wants satisfaction. Yet
retaliation has its own cost. It keeps the anger active and gives
the other person another reason to strike. Christ teaches a freedom
in which another person's hostility need not decide what we do next.
The disciple can obey him even while feeling the wrong that has
been done.

That freedom begins with the poor in spirit. Saint Augustine
understands them as humble people who submit themselves to God.
They can be taught because they accept that they need instruction.
Someone determined to justify every desire will have difficulty
hearing a command to change. Coming to Christ means allowing him
to judge what we have been calling good.

Augustine then explains why mourning can belong to conversion.
We discover that something we have loved has held us away from God.
Giving it up can hurt. The old pleasure is familiar; love for what
is truly good still needs to grow. Christ's blessing gives hope to
the person in that painful interval. Sorrow over sin can accompany
a real return to God. The desire for justice can be growing even
while obedience requires effort.

The next blessings describe that desire becoming active. The hungry
seek justice earnestly. The merciful help someone in misery.
Augustine connects the need to receive help with the willingness
to give it: the person seeking assistance from one stronger can
assist someone weaker. We come before God needing his mercy, with
other people still needing mercy from us. Humility gives us a
reason to treat them generously.

The Lesson shows us the people whose holiness we celebrate.
John hears the servants of God counted by their tribes, then sees
a multitude from every nation standing before the throne and the
Lamb. They wear white robes and hold palms. They praise God for
salvation. Their happiness has a giver, whom they acknowledge
together.

Saint Bede understands the white robes as baptism, and the palms
as victory through Christ's Cross. He also explains their praise
as the Church acknowledging what each member has received from
God. We can therefore look toward the saints with confidence
while still needing conversion ourselves. The Lord who gives
holiness also gives help for the work of living in it.

The prayers of this Mass ask for that help. We honor all the
saints and ask God for abundant mercy through their intercession.
The Sunday prayer added today asks release from sins into which
our frailty has entangled us. At the offering, another of the
Sunday prayers asks God to complete the gift he has already
granted without our merit. We are asking him to continue his
work in people who still need to learn obedience.

The Gospel then brings us to purity of heart and peace.
Augustine explains a pure heart as a single heart, directed
toward God. He compares it with an eye made fit to see light.
The promised sight of God requires the purification of the
person who will see him. Our desires matter because they can
keep pulling us away from the God we hope to know.

Peace involves those desires too. Augustine describes a person
whose impulses submit to reason, while reason itself submits
to Christ. Even the part of us that judges and directs needs
to obey him. Anger may make a demand, but Christ's command
can determine our answer. This is how peace within a person
can endure when something outside that person is painful.

Chrysostom also explains the work of peacemakers among other
people. They seek reconciliation where people are divided.
He connects their being called God's children with the Son's
work of reconciling those who have been alienated. The likeness
to God appears in what they do for others. The peace Christ
teaches extends into the relationships in which hostility
has taken hold.

Jesus immediately blesses those persecuted for justice.
Chrysostom reads the two blessings together. An agreement
requiring evil would betray the justice for which the disciple
must be willing to suffer. Jesus also specifies that the
accusations against his disciples are false and made for his
sake. Being opposed does not by itself prove that we have
acted rightly. We still have to submit our conduct to him.

This helps us understand the rest he offers. Following Christ
can require the pain of relinquishing sin and the courage to
endure another person's wrong. In that obedience, we can begin
to be free from the anger and pride which keep demanding more
from us. Chrysostom also points beyond present effort to the
glory Christ promises. The disciple's hope includes a happiness
which suffering cannot finally take away.

At Communion, the Mass repeats three of these blessings:
purity of heart, peacemaking, and persecution for justice.
Those approaching the sacrament hear again the promise of
seeing God and being called his children. We ask for the
saving benefit of our offering through God's mercy. We need
his gift as we continue learning to live the life his Son
has described.

If you have been prolonging a quarrel by returning an insult,
let your next answer be free of that retaliation. You can
address what is wrong without adding another injury. Such
restraint may leave something difficult unresolved, and the
other person remains responsible for the answer they give.
Your own answer can begin to obey Christ's command of peace.

The Offertory speaks of righteous souls in God's
hand, even when their death looks like defeat. Christ's
promise of the kingdom remains true when fidelity costs
life itself. The saints we honor have not been lost to the
God they served.

After Communion we ask for lasting joy in the saints and
the help of their prayers. We still need that help as we
learn Christ's humility and seek peace in justice. He has
promised the pure of heart that they shall see God. That
is the happiness he offers you, and the fulfillment for
which we rejoice with all his saints today.
